\begin{abstract}
\noindent
This note documents the models used in the paper \emph{The Environmental Macroeconomics of the Circular Economy}. It contains two parts: a social planner's problem, and the corresponding decentralized market economy together with the instruments that implement the first best. The quantitative implementation, with explicit functional forms, is documented in \quantnote{}, a separate document that refers to this one by section name. A defining feature of the models is that the material balance principle (mass preservation and increasing entropy) is imposed for each economic activity -- production, extraction, exploration, waste collection and treatment, consumption, investment, and recycling -- with aggregate material balances derived from the process-level accounting. Doing so ties the use of materials, through extraction of raw materials from nature and recovery of secondary materials from waste, together with waste generation and pollution, and yields an endogenous material intensity of final goods and capital.


Two additional primitives distinguish the model from the standard environmental Ramsey problem further. Waste accumulates in an anthropogenic stock that is drawn down for handling at a fixed share per period, so that circular throughput is bounded by the matter the economy has actually retained; and production carries a minimum material requirement below which it cannot run. Together they make the long-run question a survival question rather than a convergence question, and the planner part characterises the resulting taxonomy. The classification maintains trending technology, under which three long-run states are available: collapse, where production ceases after a finite operating time and consumption converges to zero along the aftermath; balanced dematerialization, where the material block decays exponentially while the goods block grows; and perpetual circular growth, where a closed loop carries a constant physical throughput forever. Whether production has a material floor and whether the recovery yield ceiling is hard or soft give four cells, and each cell gets a state: without a floor the economy survives in every case, and with a floor and a hard ceiling it collapses however fast technology improves, so that survival and perpetual circularity are the same thing and reduce to a condition on a stock, that the material the economy retains clears the floor at its own turnover rate. The stationary case, in which a closed loop is infeasible whatever the technology, is carried as a benchmark appendix outside the classification.


The market part shows that the shadow price of waste is an ordinary market price --- the gate fee of a disposal market --- whenever the waste stock is privately owned, and that three externalities remain: the pollution stock, common-pool exploration, and an attribution failure specific to the materials accounting, under which a competitive economy charges downstream users of the final good for material they did not in fact draw into circulation. Recycling requires no subsidy of its own. The quantitative note implements the planner and market economies as a single residual system with the planner nested as a corner of the policy space, and asks where a global economy calibrated from the mid-nineteenth century onward sits relative to the survival margin the theory identifies.


The appendices carry the full derivation and decomposition of the optimality conditions; the three long-run states, with their proofs and their workhorse constructions; sufficiency; the stationary benchmark; the workhorse functional forms; and a list of extensions.
\end{abstract}
